\scanpage{59}
% archive sheet number: 58
\[
\resizebox{\linewidth}{!}{$\displaystyle f\longmapsto\left(\bigl(\underbrace{f_{\deleted{\unsure{0}}}(u_x)}_{\ell_x}\bigr)_{x\in C_0},
\bigl(f_{a_i,a_i}:\underbrace{\operatorname{Aut}(a_i)\longrightarrow\deleted{\unsure{\operatorname{Aut}|)}}\operatorname{Aut}}_{\varphi_i}(f_0(a_i))\bigr)\right).$}
\]

c) Pour que $f$ soit \unsure{fidèle|fidèl}, il f et suffit que $\deleted{\unsure{\forall}}\forall i\in I$, $\varphi_i$ soit injectif.

Pour que $f$ soit pleinement \unsure{fidèle|fidèl}, il f et suffit que \deleted{\unsure{;}} 1°) pour $i\ne j$, on ait $\operatorname{Hom}(f_0(a_i),f_0(a_j))=\varnothing$

2°) pour $\deleted{\unsure{\forall i|\forall}}\ x,y\in {C_{(i)}}_0$,
\[
\deleted{\unsure{\operatorname{Hom}}}\operatorname{Isom}(f_0(x),f_0(y))\xrightarrow{\sim}\operatorname{Hom}(f_0(x),f_0(y))
\]

3°) $\forall i\in I$, $\varphi_i$ est un iso.

\deleted{d) \unsure{Soit} $C$}

Pour que $\deleted{\unsure{d}}f$ soit une équivalence, il f. et s. qu'on ait de plus

4°) Tout $x'\in C'_0$ est isomorphe à un des $f_0(a_i)$

Pour que $f$ soit un isom, il faut et s. qu'on ait 1°), 2°), 3°) et de plus

$4'$\textsuperscript{o}) $f_0$ bijective.

